%% -*- coding: utf-8 -*-

\section{180个英语俚语}

\begin{center}
  \textbf{A}
\end{center}

\begin{multicols}{2}
  \setlength{\columnseprule}{0.4pt}
  \begin{description}% [nosep]
    \item[A diamond in the rough] 有潜质的人或物；璞玉
    \item[A drop in the bucket] 微不足道；分量很少（{\fzxk 《圣经》}）
    \item[A picture is worth a thousand words] 百闻不如一见
    \item[A piece of cake] 轻而易举的事
    \item[Absence makes the heart grow fonder] 小别胜新婚（{\fzxk 莎士比亚《奥赛罗》}）
    \item[Achilles' heel] 唯一弱点，致命伤（{\fzxk 《荷马史诗》}）
    \item[All hell broke loose] 失去控制（{\fzxk 约翰·弥尔顿《失乐园》}）
    \item[All that glitters is not gold] 金玉其外败絮其中；勿以貌取人（{\fzxk 莎士比亚《威尼斯商人》，glitter被用为glister}）
    \item[An albatross around (one's) neck] 挥之不去的阴影（{\fzxk 萨缪尔·泰勒·柯立芝《老水手之歌》, 1798}）
    \item[An arm and a leg] 非常昂贵；昂贵的代价
    \item[Armed to teeth] 全副武装
    \item[At death's door] 濒临死亡（{\fzxk 《圣经·诗篇》}）
    \item[At the drop of a hat] 立刻，马上（{\fzxk 美国西部过去的决斗和赛马之类的比赛时，裁判会以扔帽子的形式宣布比赛开始。}）
    \item[At the eleventh hour] 来得早不如来得巧；关键时刻（{\fzxk 《圣经》}）
    \item[At the end of (one's) rope] 忍无可忍（{\fzxk rope最开始是用tether代替的。tether：（拴牲畜的）拴绳}）
    \item[Axe to grind] 别有用心；有不同意见（{\fzxk 本杰明·富兰克林的故事。常用法：have \(\sim\)}）
  \end{description}
\end{multicols}

\begin{center}
  \textbf{B}
\end{center}

\begin{multicols}{2}
  \setlength{\columnseprule}{0.4pt}
  \begin{description}% [nosep]
    \item[Back to square one] 前功尽弃（{\fzxk 源于英国孩童游戏，如跳房子和蛇梯棋}）
    \item[Ballpark figure] 大约的数目（{\fzxk ballpark：棒球场，始于20世纪50年代}）
    \item[Ball's in your court] 该是你采取行动的时候了（{\fzxk 来源于网球运动}）
    \item[Barking/bark up the wrong tree] 找错对象（{\fzxk 起源于狩猎}）
    \item[Beat around the bush] 拐弯抹角；旁敲侧击
    \item[Beat it] 滚开（{\fzxk it本指地面，即用脚打击地面，也就是“跑开，滚开”的意思。迈克尔·杰克逊的《Beat it》也应该翻译为“远离（暴力）”}）
    \item[Beat the rap] 全身而退（{\fzxk 源于美国的诉讼领域，指逃过刑事责任。rap通常指“惩罚”，直译即“战胜惩罚”。}）
  \end{description}
\end{multicols}

\begin{center}
  \textbf{C}
\end{center}

\begin{multicols}{2}
  \setlength{\columnseprule}{0.4pt}
  \begin{description}[style=unboxed]% [nosep]
    \item[(The) Cat's out of the bag/Let the cat out of the bag] 走漏风声（{\fzxk 源于200年前的市集，人们买肉付钱后，有些黑心卖价在最后时刻把商品拿出来装一只猫进去混淆视听，但如果猫逃出来，那就走漏风声了。}）
    \item[Catch (someone) red-handed] 逮个正着（{\fzxk 杀人犯作案后，来不及处理手上的鲜血就被抓获。}）
    \item[Chew the fat] 闲聊；八卦 （{\fzxk{}起源于19世纪，吃零食增加脂肪。}）
  \end{description}
\end{multicols}

\begin{center}
  \textbf{D}
\end{center}

\begin{multicols}{2}
  \setlength{\columnseprule}{0.4pt}
  \begin{description}[style=unboxed]% [nosep]
    \item[Dirt poor] 极其贫困的，一贫如洗 （{\fzxk{}源于16世纪，穷人家买不起木地板，家里到处是灰尘。}）
    \item[Don't count your chickens (before they are hatched)] 话别说的太早；不要打如意算盘 （{\fzxk{}《伊索寓言》}）
    \item[Don't look a gift horse in the mouth] 不要挑剔别人送的礼物 （{\fzxk{}有些人收到别人送的马，会通过看马的牙齿来判断马匹是否精良，这是不礼貌且不正确的做法。}）
    \item[Double whammy] 祸不单行 （{\fzxk{}whammy是“打击、晦气、倒霉”的意思，源于1951年漫画家艾尔·盖普的漫画《小阿布纳》。}）
    \item[Down in the dumps] 闷闷不乐；打不起精神 （{\fzxk{}down用来形容人的心情不好}）
    \item[Down to the wire] 不到最后关头难分胜负；拮据的 （{\fzxk{}源于赛马比赛，wire指跑到终点拉的终点线。}）
    \item[Draw a blank] 忘记；脑中一片空白 （{\fzxk{}源于伊丽莎白一世发型彩票，空白的则为没有中奖。}）
  \end{description}
\end{multicols}

\begin{center}
  \textbf{E}
\end{center}

\begin{multicols}{2}
  \setlength{\columnseprule}{0.4pt}
  \begin{description}[style=unboxed]% [nosep]
    \item[Eagle eye/Eagle-eyed] 眼神锐利的；有眼光的
    \item[(one's) Ears and burning] 耳朵发烫 （{\fzxk{}源于古罗马的迷信：如果别人说你坏话，你的耳朵就会发烫。}）
    \item[Eat (one's) heart out] 伤心欲绝 （{\fzxk{}《圣经》}）
    \item[Egg on (one's) face] 丢脸；尴尬 （{\fzxk{}源于20世纪60年代的美国，观众看演出时如果觉得演得不好，就会往台上砸鸡蛋。}）
    \item[Every cloud has a silver lining] 天无绝人之路；凡事都有好的一面 （{\fzxk{}英国诗人约翰·弥尔顿《酒神可莫斯》，1634。}）
    \item[Everything but the kitchen sink] 万事俱备 （{\fzxk{}厨房的洗碗柜带不走，源于20世纪初期。}）
    \item[(one's) Eyes are bigger than (one's) stomach] 贪多吃不下；眼高手低 （{\fzxk{}源于16世纪。}）
  \end{description}
\end{multicols}


\begin{center}
  \textbf{F}
\end{center}

\begin{multicols}{2}
  \setlength{\columnseprule}{0.4pt}
  \begin{description}[style=unboxed]% [nosep]
    \item[Face the music] 面对困难，承担后果 （{\fzxk{}早期歌剧院，演员面向乐池，有时候出现忘词等尴尬情况，也只能对着音乐不断的乐池干瞪眼。}）
    \item[Fair and square] 光明正大 （{\fzxk{}源于17世纪，许多木工制品都是方正的，表示没有偷工减料。}）
    \item[Fall on deaf ears] 石沉大海 （{\fzxk{}《圣经》：turn a deaf ear，表示不重视某件事。}）
    \item[(A) Far cry from] 相差很多 （{\fzxk{}原意指远到让人哭泣。}）
    \item[Feather in (one's) cap] 丰功伟业 （{\fzxk{}源于印第安人满头的羽毛。}）
    \item[Fight fire with fire] 以火攻火；以毒攻毒 （{\fzxk{}源于19世纪中期，美国西部的人民为了防止敌人防火，先把一块空地烧光作为防火带。}）
    \item[Fly-by-night] 不可信任的，不可靠的 （{\fzxk{}源于以前的巫婆总是在夜里骑着扫帚飞过，而那时巫婆就是骗子的代名词。}）
    \item[Fly off the handle] 大发雷霆 （{\fzxk{}据说源于美国西部拓荒时期，人们用斧头砍树，但是斧头质量有时太差，斧柄经常掉落，人们就会大发雷霆。}）
    \item[For the bird] 毫无价值的；糟糕的 （{\fzxk{}据说以前马车长途前进时，休息时马拉完粪便后，会有小鸟过来啄食。}）
    \item[Foregone conclusion] 不可避免的结局 （{\fzxk{}莎士比亚《奥赛罗》}）
    \item[(A) Frog in (one's) throat] 喉咙不舒服 （{\fzxk{}据说以前人们在河里饮水不小心喝下青蛙卵，人们相信卵会在喉咙里变成青蛙，哽住喉咙。}）
  \end{description}
\end{multicols}

\begin{center}
  \textbf{G}
\end{center}

\begin{multicols}{2}
  \setlength{\columnseprule}{0.4pt}
  \begin{description}[style=unboxed]% [nosep]
    \item[Get the sack] 被开除 （{\fzxk{}过去人们上班会带一个大布袋，里面装上自己工作会用到的工具等，老板要辞退他，就会把大袋子丢还给他，意即让他走人。}）
    \item[Gift of the gab] 能言善道 （{\fzxk{}gab指“空谈”，自苏格兰文中gabba演化而来。}）
    \item[Give me a break] 饶了我吧
    \item[Give the thumbs up/down] 称赞/贬抑
    \item[Go berserk] 发狂 （{\fzxk{}berserk源自北欧古语言，指“披着熊皮的人”，后来引申为“狂战士”，好战且毫无恐惧和疼痛感，进入战斗状态后会陷入癫狂状态。}）
    \item[Go Dutch] 各自付账； AA （{\fzxk{}17史记荷兰商人就已经有了各自付费的习惯，因为商人流动性很大，也许再也见不到面了。}）
    \item[Go out on a limb] 冒险 （{\fzxk{}limb指树枝，现在已经引申为不单单指身体上的冒险，更多的指有辱名誉、丢脸或被误会的风险等。}）
    \item[Goody two-shoes] 假正经；伪善者 （{\fzxk{}源自童话《The history of Little Goody Two-Shoes》，原来的意思是是如果你是一个好心又正直的人，就会遇到很多好事，现在已经演变为了贬义词。}）
    \item[Green with envy] 嫉妒；羡慕 （{\fzxk{}《奥赛罗》中使用 green eyes形容人的嫉妒。}）
    \item[Gung ho] 充满热情的；带劲的 （{\fzxk{}二战时一位美国上校被派到中国，这个词来自中文“共和”，表示合作共赢，后来演变为他激励将士的话语。}）
    \item[Gut feeling] 直觉；预感 （{\fzxk{}《圣经》}）
  \end{description}
\end{multicols}

\begin{center}
  \textbf{H}
\end{center}

\begin{multicols}{2}
  \setlength{\columnseprule}{0.4pt}
  \begin{description}[style=unboxed]% [nosep]
    \item[Hair-raising] 令人毛骨悚然的
    \item[Hand over fist] 快速 （{\fzxk{}主要形容赚钱很快。或源于航海，水手需要手动拉绳索来摆动船帆，动作方式是：一只手抓住绳子，握成拳 fist， 另一只手迅速往上再抓，两只手交替进行。}）
    \item[Hang out] 和朋友一起打发时间；一起出去玩
    \item[Hard up] 生活拮据 （{\fzxk{}源于航海，在风暴到来时，水手应当调转船头，冲着逆风的方向，这样比较容易改变方向。}）
    \item[Hit the nail on the head] 一针见血；切中问题要害 （{\fzxk{}此处的head指钉头，只有砸在钉头，才能钉得最深。}）
    \item[Hit the sack] 就寝 （{\fzxk{}以前的被单床罩都是用麻布做的。}）
  \end{description}
\end{multicols}

\begin{center}
  \textbf{I}
\end{center}

\begin{multicols}{2}
  \setlength{\columnseprule}{0.4pt}
  \begin{description}[style=unboxed]% [nosep]
    \item[Icing on the cake] 锦上添花 （{\fzxk{}原意指“蛋糕上的糖霜”。}）
    \item[If the show fits (wear it)] 接受惩罚；放手去做 （{\fzxk{}shoe源于早期学生受惩罚时戴的纸帽。}）
    \item[Ignorance is bliss] 无知就是福 （{\fzxk{}托马斯·格雷的文章，1742。bliss指“极乐”。}）
    \item[In a nutshell] 总而言之；长话短说
    \item[In the black/red] 获利/亏损 （{\fzxk{}会计在做账时，盈利的项目用黑笔标识，亏损的用红笔。}）
    \item[In the doghouse] 惹上麻烦；成为别人生气的对象 （{\fzxk{}据说有些家长在孩子调皮时，吓唬他们要把他们关到外面的狗屋里。}）
    \item[In the nick of time] 在最后一刻；紧要关头 （{\fzxk{}nick是木头上的刻痕，以前人们计时会在木头上刻痕。}）
    \item[It's not over till the fat lady sings] 别妄下结论；不到最后不见真章 （{\fzxk{}歌剧的结束前的高潮往往是一位身材丰满的女声乐家的表演。}）
  \end{description}
\end{multicols}

\begin{center}
  \textbf{J}
\end{center}

\begin{multicols}{2}
  \setlength{\columnseprule}{0.4pt}
  \begin{description}[style=unboxed]% [nosep]
    \item[Jack of all trades (master of none)] 样样都懂，却样样都不精通 （{\fzxk{}源自17世纪英国，最初用来形容某人十八般武艺样样精通，后面加上master of none，就变成了负面的意思。}）
    \item[Jekyll and Hyde] 双重人格 （{\fzxk{}英国·史蒂文森《化身博士》，主角Jekyll本是善良的博士，将自己当作实验对象，结果却导致人格分裂，变成邪恶的Hyde。}）
    \item[(The) Jig is up] 到此为止 （{\fzxk{}jig原来形容一种技巧或者骗术，常被用在假幽默或者耍宝的人身上。}）
    \item[Jump the gun] 操之过急 （{\fzxk{}源于田径赛跑之前都会有枪发令，运动员才开始起跑，否则就是抢跑。}）
    \item[Just around the corner] 即将到来
  \end{description}
\end{multicols}

\begin{center}
  \textbf{K}
\end{center}

\begin{multicols}{2}
  \setlength{\columnseprule}{0.4pt}
  \begin{description}[style=unboxed]% [nosep]
    \item[Keep it under your hat] 保守秘密 （{\fzxk{}源于中古时期，弓箭手在打仗时，会将弓弦藏在帽子里以备不时之需，同时保持其干燥。后来人们在紧急情况突发时，会将重要的东西藏在帽子里。}）
    \item[Keep your fingers crossed] 祈求好运 （{\fzxk{}手指交错，代表十字架，也是基督教徒相认的一种方式。}）
    \item[Keep your shirt on] 别太激动；别生气 （{\fzxk{}19世纪的人们没钱买昂贵的衣服，因此想和别人动手前，先把衣服脱掉以防弄坏。}）
    \item[Kick the bucket] 死亡；翘辫子 （{\fzxk{}起源一来自屠宰场的猪，屠夫将猪绑在一块木板上，猪知道死亡将近，挣扎而脚不断踢那块木板，即bucket；起源二则是形容人上吊自杀时，一脚踢开脚下的板凳，即bucket。}）
    \item[Kick two birds with one stone] 一箭双雕；一石二鸟 （{\fzxk{}源自2000多年前的拉丁文。}）
    \item[Knock on wood] 老天保佑；好险！ （{\fzxk{}源自19世纪，人们迷信的认为木头中都藏着小精灵，所以在木头上敲一下可能带来好运。}）
    \item[Knuckle down] 屈服；认真工作 （{\fzxk{}源于18世纪的弹珠游戏，根据规则，双方必须把指关节，即knuckle贴在地上弹弹珠。}）
    \item[Kowtow to (someone)] 拍马屁；卑躬屈膝 （{\fzxk{}kowtow即中文的“磕头”。}）
  \end{description}
\end{multicols}

\begin{center}
  \textbf{L}
\end{center}

\begin{multicols}{2}
  \setlength{\columnseprule}{0.4pt}
  \begin{description}[style=unboxed]% [nosep]
    \item[Last straw] 最后一根稻草 （{\fzxk{}该成语是“the straw that breaks the camel's back”的简化。}）
    \item[(The) Law of the jungle] 丛林法则；弱肉强食 （{\fzxk{}童话作家卢德亚·奇普林《丛林怪谈》}）
    \item[Leave no stone unturned] 全力以赴 （{\fzxk{}古希腊传说中，有一位先知告诉寻宝人要把地上所有的大石块都翻过来才有可能找到真正的宝藏。}）
    \item[Let one's hair down] 放轻松；别拘束 （{\fzxk{}西方女子在外出或者参加宴会等正式场合都会把头发盘起来，只有到家后才会把头发放下来。}）
    \item[Lip service] 随口说说；光说不练 （{\fzxk{}《圣经》}）
    \item[(One's) Lips are sealed] 守口如瓶；保密
    \item[Live from hand to mouth] 勉强度日 （{\fzxk{}16世纪英国发生了大饥荒，有食物的时候往往饥不择食，直接把手里的事物送到嘴里，以防落入他人口中。}）
    \item[Long in the tooth] 年长的；老相的 （{\fzxk{}以前人们会用马的牙齿长短来判断其年龄。}）
    \item[Long shot] 机会渺茫 （{\fzxk{}源自19世纪，与用枪有关，目标越远越难命中。}）
    \item[Long time no see] 好久不见 （{\fzxk{}源于 Chinglish。}）
    \item[Loose cannon] 我行我素，不顾后果的人 （{\fzxk{}源自19世纪航海领域，Cannon如果没有固定好，会在甲板上滚来滚去，十分危险。}）
    \item[Lion's share] 最大、最好的部分 （{\fzxk{}《伊索寓言》狮子率众捕猎后分为四份，自己认为理所当然应该全部占有。}）
  \end{description}
\end{multicols}

\begin{center}
  \textbf{M}
\end{center}

\begin{multicols}{2}
  \setlength{\columnseprule}{0.4pt}
  \begin{description}[style=unboxed]% [nosep]
    \item[Mad as a hatter] 非常疯狂的 （{\fzxk{}起源于苏格兰帽匠，据说当时制作毡帽的皮得经过含汞的硝酸盐进行处理在加工，因此制作过程中，帽匠会吸入汞而中毒，症状为口齿不清，全身抽搐，和疯子症状类似。}）
    \item[Make a mountain out of a molehill] 小题大做 （{\fzxk{}molehill指“鼹丘”，即鼹鼠挖洞而堆积土壤成为的小土包。}）
    \item[Make (one's) hair stand on end] 恐怖吓人；让人毛骨悚然 （{\fzxk{}莎士比亚《哈姆雷特》，相反的可以用“not turn a hair”表示“毫不慌张，表现镇定自若”。}）
    \item[Mickey mouse] 低级的；品质差的 （{\fzxk{}20世纪30年代，迪士尼授权给米奇老鼠制作的文具、手表等，早期品质比较差。}）
    \item[Money is the root of all evil] 前十万恶之源 （{\fzxk{}《圣经》}）
    \item[Monkey business] 胡闹，恶作剧；骗人的把戏
    \item[Not giving a monkey's (about something)] 漠不关心
    \item[Mumbo jumbo] 莫名其妙的话；无意义的活动 （{\fzxk{}Mumbo jumbo是非洲某些地区崇拜的神的名字，祭司在祷告时往往喃喃呓语。}）
  \end{description}
\end{multicols}

\begin{center}
  \textbf{N}
\end{center}

\begin{multicols}{2}
  \setlength{\columnseprule}{0.4pt}
  \begin{description}[style=unboxed]% [nosep]
    \item[Neck and neck] 并驾齐驱；不分上下 （{\fzxk{}赛马场上，两匹马跑的差不多的时候，仅凭颈部能看出谁在前。}）
    \item[Needle in a haystack] 大海捞针 （{\fzxk{}haystack指干草堆。}）
    \item[Nest egg] 老本；存款 （{\fzxk{}原指放在巢中以引诱鸟继续在巢中产蛋的人工蛋或真蛋。}）
    \item[Nip in the bud] 防患于未然；防微杜渐 （{\fzxk{}nip指“掐”，bud指“花蕾”。}）
    \item[No holds barred] 没有限制；火力全开 （{\fzxk{}hold在摔跤术语中，hold是一种擒抱技术。}）
    \item[No love lost] 没有好感；仇恨；敌意
    \item[No strings attached] 无附带条件 （{\fzxk{}string是线、细绳，可用来控制木偶表演。}）
    \item[Nose to the grindstone] 埋头苦干
    \item[Pay through the nose] 付出比真正价值高得多的钱
    \item[Nose around/about] 打听别人的事
    \item[no skin of my nose] 不关我的事
    \item[Not all it's cracked up to be] 名不副实 （{\fzxk{}14世纪时，cracked指“充满自信的，思维敏捷的说话”。}）
    \item[Not cricket] 不公平 （{\fzxk{}板球游戏一向被称为“绅士的游戏”，崇尚体育精神和公平比赛。}）
    \item[Not up to scratch] 未达标准；不够资格 （{\fzxk{}源于拳击比赛，若一方被击倒，又不能于38秒内回到分界线继续，则另一方获胜。分界线就是scratch。}）
  \end{description}
\end{multicols}

\begin{center}
  \textbf{O}
\end{center}

\begin{multicols}{2}
  \setlength{\columnseprule}{0.4pt}
  \begin{description}[style=unboxed]% [nosep]
    \item[Off the beaten track] 不落俗套；独辟蹊径 （{\fzxk{}所谓“beaten track”字面意思是“被打的小路”，引申为走过很多次的小路，这条成语也即“不走寻常路”之意。}）
    \item[Old wives' tale] 无稽之谈 （{\fzxk{}wives此处指“生活在古代、迷信的人”。该成语可能来自古希腊先哲柏拉图。}）
    \item[On cloud nine] 快乐得不得了 （{\fzxk{}cloud nine指蓬松的积雨云，源自20世纪50年代美国气象局对云的分类。又或源于佛教的“九重云霄”之说。}）
    \item[On the ball] 机灵，反应快 （{\fzxk{}源于棒球运动，若一名投手出手不凡，成功地为对面的接球手制造了不少麻烦的话，称之为“have a lot on the ball”，因为投手会做出旋转动作或其他的一些假动作。}）
    \item[Once in a blue moon] 千载难逢 （{\fzxk{}“blue moon”非常罕见，只有在火山爆发后，空气中粉尘的大小正好能衍射出蓝色的光，使月亮看起来是蓝色的。}）
    \item[Out of the blue] 突然 （{\fzxk{}可能来源于“a bolt from blue”，即晴天霹雳。}）
    \item[Over the hill] 走下坡路；过了巅峰期
  \end{description}
\end{multicols}

\begin{center}
  \textbf{P}
\end{center}

\begin{multicols}{2}
  \setlength{\columnseprule}{0.4pt}
  \begin{description}[style=unboxed]% [nosep]
    \item[Pan out] 结果是；取得成功 （{\fzxk{}19世纪的加州是淘金圣地，淘金者用盘子舀起泥沙细细晃动，泥沙被水冲走后盘子上如果留有金子，就意味着发财了。}）
    \item[Pandora's box] 罪恶；麻烦的根源 （{\fzxk{}古希腊神话。}）
    \item[Par for the course] 果不其然；意料中的事 （{\fzxk{}par来自拉丁文，相当于“equal”，该短语源自高尔夫运动，标准杆(par)是指选手在没有犯错的情况下打完每个洞所应有的杆数。又比如：above par指水准之上，below par指水准之下。}）
    \item[Pass the buck] 推卸责任 （{\fzxk{}杜鲁门总统的名言：“the buck stops here”，指“他会为政府部门的行为承担最终的责任”，buck意指“责任”。}）
    \item[Pass with flying colors] 高分过关 （{\fzxk{}17世纪曾有短语“to come off with flying colors”，指那些得胜归来的战船。}）
    \item[Pay through the nose] 付出大笔钱 （{\fzxk{}9世纪时，爱尔兰被丹麦统治，被迫缴纳高额税金，而逃税被逮住的人，常常被割掉鼻子以示惩罚。}）
    \item[Pipe dream] 白日梦，空想 （{\fzxk{}指19世纪人们用烟斗吸食鸦片，描述他们吞云吐雾的状态。}）
    \item[Play ball] 合作 （{\fzxk{}“You play ball with me and I'll play ball with you.”指你和我打球，我以后也会和你打球。}）
    \item[Play (it) by ear] 即兴，随性的 （{\fzxk{}源于音乐领域，一些乐手不懂五线谱，只能通过耳朵听，然后模仿演奏，也就是即兴表演。}）
    \item[Play (one's) chain] 愚弄人；惹人生气 （{\fzxk{}以前猎手抓到猛兽，会用链子把它们绑住，一旦拉扯链子，会使野兽勃然大怒。}）
    \item[Pull (one's) leg] 开玩笑；逗弄别人 （{\fzxk{}引申自“扯某人后腿，害人跌倒出糗”之意。}）
    \item[Pull out all the stops] 全力以赴 （{\fzxk{}stop指老式管风琴的音栓，当风琴手把所有音栓都拉出，那么众管齐鸣，发出最大的声音。}）
    \item[Pull the wool over (one's) eyes] 欺骗；蒙蔽 （{\fzxk{}起源于18世纪流行的男士假发，特别是法官和律师戴的那种带卷的假发。Wool则是假发的俗称。}）
    \item[Push the envelope] 挑战极限 （{\fzxk{}在航空领域，对飞行员来说，envelope指安全范围内操控飞机的极限，飞行员“push the envelope”指要把飞机的性能发挥到极限。}）
  \end{description}
\end{multicols}

\begin{center}
  \textbf{R}
\end{center}

\begin{multicols}{2}
  \setlength{\columnseprule}{0.4pt}
  \begin{description}[style=unboxed]% [nosep]
    \item[Rack (one's) brain] 绞尽脑汁 （{\fzxk{}rack指“折磨；榨取”。}）
    \item[Rain cats and dogs] 倾盆大雨 （{\fzxk{}源自伦敦一场大雨之后淹死了很多猫和狗。}）
    \item[Rain check] 延期 （{\fzxk{}源于19世纪80年代，当时棒球运动大受欢迎，有很多露天举行的棒球比赛，但要是碰到下雨，就要暂停延期，观众可以领到“雨票”，或用原票根作为“雨票”，可以在改期后凭其入场。}）
    \item[Rat race] 激烈的竞争；无谓的奔忙 （{\fzxk{}想象一群饥饿的老鼠看见好吃的一拥而上的情形吧。}）
    \item[Read between the lines] 弦外之音
    \item[Ring the bell] 想起；忽然有了模糊的记忆 （{\fzxk{}早期教堂有重要事情都会敲钟传达信息，人们也可以根据钟声集合以便应付突发事件。}）
    \item[Rip off/Rip-off] 敲竹杠；冤大头 （{\fzxk{}前者为动词，后者为名词。}）
    \item[Risk life and limb] 冒着生命危险 （{\fzxk{}limb指四肢。}）
    \item[Rule of thumb] 经验法则 （{\fzxk{}起源一来自以前人们习惯用大拇指的中间指节到指甲尖作为天然的测量工具，大约一英寸；起源二说有经验的酿酒师习惯用大拇指伸进酿酒池测量发酵过程的温度，以了解酿造进展。}）
    \item[Run out of steam] 筋疲力尽 （{\fzxk{}源于蒸汽时代，蒸汽越少动力越弱。}）
  \end{description}
\end{multicols}

\begin{center}
  \textbf{S}
\end{center}

\begin{multicols}{2}
  \setlength{\columnseprule}{0.4pt}
  \begin{description}[style=unboxed]% [nosep]
    \item[Save a face] 保全颜面 （{\fzxk{}“lose face”则指“丢脸”。}）
    \item[Seal (one's) fate] 决定（某人的）命运
    \item[See the light] 恍然大悟
    \item[Settle a/the score] 讨公道；算旧账；报复 （{\fzxk{}settle表示“支付、结算”，score表示“【旧时酒馆等的】帐，欠账”。}）
    \item[Sitting duck] 坐以待毙；待宰羔羊 （{\fzxk{}源于打猎，比起飞在天上的野鸭，坐着不动的是更容易命中的目标。}）
    \item[Smart aleck/alec] 自作聪明的人 （{\fzxk{}埃里克·侯格【Alec Hoag】是19世纪纽约知名的小偷，其欺骗人的技巧十分高名。}）
    \item[Start from scratch] 白手起家 （{\fzxk{}19世纪的体育运动用语，运动员在赛跑前，都会用棍子在起点画一道横线，这道线就是“scratch”。}）
    \item[Start the ball rolling] 开始进行 （{\fzxk{}起源于19世纪的槌球运动，用木槌击球，使其穿过金属环，要开始比赛，就要先“让球滚动”。}）
    \item[Strike while the iron is hot] 趁热打铁
  \end{description}
\end{multicols}

\begin{center}
  \textbf{T}
\end{center}

\begin{multicols}{2}
  \setlength{\columnseprule}{0.4pt}
  \begin{description}[style=unboxed]% [nosep]
    \item[Take a back seat] 退居二线
    \item[Take pains] 煞费苦心；费尽心力 （{\fzxk{}源自一个古老的短语“pain yourself”，指“特别照顾”或“努力”。}）
    \item[Take to the cleaners] 洗劫一空；骗光（某人的）钱
    \item[That's the ticket] 这就对了；这正是所需的 （{\fzxk{}ticket或许来自法语的“etiquette”，指“得体的行为”。}）
    \item[Think outside the box] 打破常规；跳出传统思考模式
    \item[Through thick and thin] 在任何情况下；同甘共苦 （{\fzxk{}来源于中世纪英国的森林环境：人们在森林中赶路，会经历两种地形：thick = 茂密丛林（难走）；thin = 稀疏树林（好走）。所以“through thick and thin” 本来就是“穿过密林和疏林”。}）
    \item[Two heads are better than one] 三个臭皮匠顶个诸葛亮
  \end{description}
\end{multicols}

\begin{center}
  \textbf{U\(-\)W}
\end{center}

\begin{multicols}{2}
  \setlength{\columnseprule}{0.4pt}
  \begin{description}[style=unboxed]% [nosep]
    \item[Under the gun] 压力之下
    \item[Under the weather] 不舒服 （{\fzxk{}在过去的航海时代，海上经常有恶劣天气，船员容易晕船、生病，生病的人会被安排到甲板下（below deck）休息，也就是“在天气之下”（under the weather）。}）
    \item[Up in arms] 怒气冲冲 （{\fzxk{}原意指“拿起武器准备战斗”。}）
    \item[Vested interests] 既得利益 （{\fzxk{}心理学教授威廉·克拉诺于1995年提出，这是一种心理现象，是结果影响先行行为的典型。形容如果已经尝到甜头，即使知道这件事不对也会坚决支持。}）
    \item[Wear (one's) heart on (one's) sleeve] 流露感情；喜怒形于色 （{\fzxk{}源自中世纪的格斗比赛，骑士会在胳膊上系上丝带，而丝带的颜色代表支持他的女性。}）
    \item[Wet behind the ears] 缺乏经验的人；乳臭未干 （{\fzxk{}大多数胎生动物的幼崽出生后，浑身都是湿的，身上渐渐干了之后，耳朵后面还是湿的。}）
    \item[When pigs fiy] 不可能发生的事 （{\fzxk{}《爱丽丝梦游仙境》}）
    \item[(A) Wild goose chase] 徒劳之举 （{\fzxk{}源于赛马活动，赛马时，马群会跟着头马奔跑，就像野鹅的飞行阵型。人是不可能追上飞行的野鹅的。}）
    \item[(The) Writing is on the wall] 不祥之兆；大难临头 （{\fzxk{}来自《圣经·但以理书》：巴比伦国王在宴会中狂欢，突然出现一只“神秘的手”，在墙上写下一行字，内容是“Mene, Mene, Tekel, Upharsin”，含义（被先知 Daniel 解读）为：“王国要灭亡，审判已经到来。”结果当天晚上王被杀，国家灭亡。}）
  \end{description}
\end{multicols}

\begin{center}
  \textbf{Y}
\end{center}

\begin{multicols}{2}
  \setlength{\columnseprule}{0.4pt}
  \begin{description}[style=unboxed]% [nosep]
    \item[You are what you eat] 吃什么像什么
    \item[You can't have your cake and eat it too.] 鱼与熊掌不可兼得 （{\fzxk{}16世纪英国剧作家约翰·海伍德创作。}）
    \item[You can't teach an old dog new tricks] 老狗变不出新把戏；人老则守旧
  \end{description}
\end{multicols}




%%% Local Variables:
%%% TeX-engine: xetex
%%% TeX-master: "../IELTS_300_Sentences"
%%% End:

% LocalWords:  vlines hlines halign valign xetex LocalWords nosep utf Chinglish
